\chapter{1959年全国卷}

\section{解答题}

\begin{problem}
\begin{enumerate}
\item 已知 \(\lg 2=0.3010\)，\(\lg 7=0.8451\)，求 \(\lg 35\).
（注）\(\lg\) 是以 \(10\) 为底的对数的符号.

\item 化简 \(\frac{\left( 1-\i \right)^3}{1+\i}\).

\item 解不等式：\(2x^2+5x<3\).

\item 求 \(\cos 165^\circ\) 的值.

\item 一直圆台上底的面积为 \(25\pi \mathrm{~cm}^2\)，下底的直径为 \(20 \mathrm{~cm}\)，母线长为 \(10 \mathrm{~cm}\). 求这直圆台的侧面积.

\item 有三条不在同一平面内的平行线 \(a\)，\(b\) 和 \(c\). 在线 \(a\) 上取一固定线段 \(AB\)，在线 \(c\)，\(b\) 上各任取一点 \(C\) 和 \(D\). 求证：不论 \(C\) 和 \(D\) 取在 \(c\)，\(b\) 的什么位置上，四面体 \(ABCD\) 的体积总是不变的.

\end{enumerate}
\bitmapfigure[max width=0.34\linewidth,trim=0 0.35cm 0 0.1cm,clip]{img/1959/national/q1_fig1.png}
\end{problem}

\begin{solution}
\begin{enumerate}
\item 因为 \(35=\frac{10\times 7}{2}\)，所以由对数的积、商公式，
\[
\lg 35=\lg 10+\lg 7-\lg 2=1+0.8451-0.3010=1.5441
\]

\item 因为 \(\left(1+\i\right)\left(1-\i\right)=2\)，所以
\[
\frac{\left(1-\i\right)^3}{1+\i}=\frac{\left(1-\i\right)^4}{\left(1+\i\right)\left(1-\i\right)}=\frac{\left(1-\i\right)^4}{2}
\]
又因为 \(\i^2=-1\)，从而 \(\i^3=-\i\)、\(\i^4=1\)，故
\[
\frac{\left(1-\i\right)^3}{1+\i}=\frac{1-4\i+6\i^2-4\i^3+\i^4}{2}=\frac{1-4\i-6+4\i+1}{2}=-2
\]

\item 移项，得
\[
2x^2+5x-3<0
\]
因式分解得 \(\left(2x-1\right)\left(x+3\right)<0\). 两个因式的零点分别是 \(x=\frac12\) 和 \(x=-3\). 当 \(x<-3\) 时两个因式均为负，乘积为正；当 \(-3<x<\frac12\) 时 \(2x-1<0\)、\(x+3>0\)，乘积为负；当 \(x>\frac12\) 时两个因式均为正，乘积为正. 因此原不等式的解为
\[
-3<x<\frac12
\]

\item 利用余弦的差角公式，得
\[
\begin{gathered}
\cos 165^\circ=\cos\left(180^\circ-15^\circ\right)=-\cos15^\circ=-\cos\left(45^\circ-30^\circ\right)\\
=-\left(\cos45^\circ\cos30^\circ+\sin45^\circ\sin30^\circ\right)=-\left(\frac{1}{\sqrt2}\cdot\frac{\sqrt3}{2}+\frac{1}{\sqrt2}\cdot\frac12\right)\\
=-\frac{\sqrt6+\sqrt2}{4}
\end{gathered}
\]

\item 设直圆台上、下底的半径分别为 \(r\)、\(R\)，母线长为 \(L\). 由上底面积，
\[
\pi r^2=25\pi
\]
所以 \(r=5\mathrm{~cm}\). 由下底直径为 \(20\mathrm{~cm}\)，得 \(R=10\mathrm{~cm}\)，且 \(L=10\mathrm{~cm}\).

将圆台侧面展开为扇环，设扇环的外、内半径分别为 \(l_1\)、\(l_2\)，圆心角为 \(\alpha\)（用弧度表示）. 外弧和内弧分别是下、上底的圆周，所以
\[
\begin{cases}
\alpha l_1=2\pi R\\
\alpha l_2=2\pi r
\end{cases}
\]
又因 \(l_1-l_2=L\)，两式相减得 \(\alpha L=2\pi\left(R-r\right)\). 因而扇环面积为
\[
S=\frac12\alpha\left(l_1^2-l_2^2\right)=\frac12\alpha\left(l_1-l_2\right)\left(l_1+l_2\right)
\]
由 \(l_1:l_2=R:r\) 及 \(l_1-l_2=L\)，可得 \(l_1+l_2=\frac{L\left(R+r\right)}{R-r}\)，所以
\[
S=\frac12\cdot\frac{2\pi\left(R-r\right)}{L}\cdot L\cdot\frac{L\left(R+r\right)}{R-r}=\pi L\left(R+r\right)
\]
因而
\[
S=\pi\cdot10\cdot\left(10+5\right)=150\pi\mathrm{~cm}^2
\]

\item 设直线 \(a\) 和 \(b\) 所确定的平面为 \(\Pi\). 由于 \(a\parallel b\)，两条平行线间的距离为定值，记为 \(d\). 对任意取在 \(b\) 上的点 \(D\)，以 \(AB\) 为底边时，三角形 \(ABD\) 的高就是点 \(D\) 到直线 \(a\) 的距离，因此
\[
S_{\triangle ABD}=\frac12 AB\cdot d
\]
这个面积与点 \(D\) 在直线 \(b\) 上的位置无关.

又因为 \(c\parallel a\)、\(c\parallel b\)，且三条直线不在同一平面内，所以直线 \(c\) 平行于平面 \(\Pi\)，而不在平面 \(\Pi\) 内. 因此，直线 \(c\) 上任意一点 \(C\) 到平面 \(\Pi\) 的距离都是同一个定值，记为 \(h\). 由于 \(A\)、\(B\)、\(D\) 都在 \(\Pi\) 内，平面 \(ABD\) 就是 \(\Pi\)，所以四面体 \(ABCD\) 的高为 \(h\). 从而
\[
V_{ABCD}=\frac13S_{\triangle ABD}h=\frac16AB\cdot d\cdot h
\]
右端只含固定的 \(AB\)、\(d\) 和 \(h\)，故无论 \(C\)、\(D\) 在各自直线上的什么位置，四面体 \(ABCD\) 的体积都不变.
\end{enumerate}
\end{solution}

\begin{problem}
三个数成等差数列，前两个数的和的三倍等于第三个数的二倍；如果第二个数减去 \(2\)（仍当作第二项），三个数就成等比数列. 求原来的三个数.
\end{problem}

\begin{solution}
设原来的三个数依次为 \(x-y\)、\(x\)、\(x+y\)，其中 \(y\) 是公差. 由第一个条件，
\[
3\left[\left(x-y\right)+x\right]=2\left(x+y\right)
\]
化简得
\[
4x=5y
\]
第二个数减去 \(2\) 后，三个数为 \(x-y\)、\(x-2\)、\(x+y\). 三个数成等比数列时，中间项的平方等于两端项的乘积，所以
\[
\left(x-2\right)^2=\left(x-y\right)\left(x+y\right)
\]
展开并消去 \(x^2\)，得
\[
y^2-4x+4=0
\]
由前式得 \(x=\frac54y\)，代入后式，得到
\[
y^2-5y+4=0
\]
即 \(\left(y-1\right)\left(y-4\right)=0\).

当 \(y=1\) 时，\(x=\frac54\)，于是原来的三个数分别为 \(x-y=\frac14\)、\(x=\frac54\)、\(x+y=\frac94\).
此时第二个数减去 \(2\) 后为 \(-\frac34\)，三个数 \(\frac14\)、\(-\frac34\)、\(\frac94\) 的公比都是 \(-3\)，满足条件.

当 \(y=4\) 时，\(x=5\)，于是原来的三个数分别为 \(x-y=1\)、\(x=5\)、\(x+y=9\).
此时第二个数减去 \(2\) 后为 \(3\)，三个数 \(1\)、\(3\)、\(9\) 的公比都是 \(3\)，也满足条件. 所以所求的三个数为 \(\frac14\)，\(\frac54\)，\(\frac94\)，或 \(1\)，\(5\)，\(9\).
\end{solution}

\begin{problem}
设有 \(\triangle ABC\)，已知 \(\angle B=60^\circ\)，\(AC\) 边长为 \(4\)，面积为 \(\sqrt3\). 求 \(AB\) 及 \(BC\) 两边之长.
\end{problem}

\begin{solution}
设 \(AB=x>0\)、\(BC=y>0\). 在三角形 \(ABC\) 中，边 \(AC\) 对着角 \(B\)，由余弦定理，
\[
4^2=x^2+y^2-2xy\cos60^\circ=x^2+y^2-xy
\]
由两边及其夹角求面积，
\[
\sqrt3=\frac12xy\sin60^\circ=\frac12xy\cdot\frac{\sqrt3}{2}
\]
所以
\[
xy=4
\]
由前两式得 \(x^2+y^2=20\). 因而
\[
(x+y)^2=x^2+y^2+2xy=20+8=28
\]
由于 \(x\)，\(y>0\)，所以 \(x+y=2\sqrt7\). 同样，
\[
(x-y)^2=x^2+y^2-2xy=20-8=12
\]
故 \(x-y=\pm2\sqrt3\). 分别相加、相减，得到 \(x=\sqrt7\pm\sqrt3\)，\(y=\sqrt7\mp\sqrt3\).
因为 \(\sqrt7>\sqrt3\)，上述两种取值都给出正的边长. 因此 \(AB=\sqrt7+\sqrt3\)，\(BC=\sqrt7-\sqrt3\)，或 \(AB=\sqrt7-\sqrt3\)，\(BC=\sqrt7+\sqrt3\). 还需核验三角形条件：\(x+y=2\sqrt7>4=AC\)，且 \(\left\lvert x-y\right\rvert=2\sqrt3<4\)，故上述两种取值均确实构成三角形.
\FigureLayoutDeclare{img/1959/national/q3_solution_fig1.png}{4.0cm}{0bp 1bp 0bp 0bp}
\solutionfigure{\bitmapfigure[max width=0.34\linewidth]{img/1959/national/q3_solution_fig1.png}}
\end{solution}

\begin{problem}
\(A\)，\(B\)，\(C\) 是直线 \(L\) 上三点，\(P\) 是这直线外一点. 已知 \(AB=BC=a\)，\(\angle APB=90^\circ\)，\(\angle BPC=45^\circ\). 试求：

\begin{enumerate}
\item \(\angle PBA\) 的正弦，余弦，正切；
\item 线段 \(PB\) 的长；
\item \(P\) 点到直线 \(L\) 的距离.
\end{enumerate}
\end{problem}

\begin{solution}
\begin{enumerate}
\item 设 \(\theta=\angle PBA\)，并记 \(PB=x\). 因为 \(\angle APB=90^\circ\)，所以三角形 \(APB\) 在 \(P\) 处为直角，\(AB=a\) 是斜边，从而 \(x=PB=a\cos\theta\).
由于 \(A\)、\(B\)、\(C\) 是直线上的三个不同点且 \(AB=BC\)，所以 \(B\) 是 \(AC\) 的中点，射线 \(BA\) 与射线 \(BC\) 反向，所以
\(\angle PBC=180^\circ-\theta\). 在三角形 \(PBC\) 中，已知 \(\angle BPC=45^\circ\)，故 \(\angle PCB=180^\circ-\left(180^\circ-\theta\right)-45^\circ=\theta-45^\circ\). 正弦定理给出 \(\frac{a}{\sin45^\circ}=\frac{x}{\sin\left(\theta-45^\circ\right)}\). 将前式代入，并利用 \(\sin45^\circ=\cos45^\circ\)，得 \(\sin\left(\theta-45^\circ\right)=\cos\theta\sin45^\circ\)，即 \(\sin\theta\cos45^\circ-\cos\theta\sin45^\circ=\cos\theta\sin45^\circ\). 因此 \(\sin\theta=2\cos\theta\)，即 \(\tan\theta=2\). 角 \(\theta\) 为锐角，所以 \(\sin\theta=\frac{2}{\sqrt5}\)，\(\cos\theta=\frac{1}{\sqrt5}\)，\(\tan\theta=2\).

\item 由前式，
\(PB=x=a\cos\theta=\frac{a}{\sqrt5}\).
\item 设从 \(P\) 向直线 \(L\) 作垂线，垂足为 \(H\)，则 \(PH\) 是所求距离. 在直角三角形 \(PBH\) 中，
\(PH=PB\sin\theta=\frac{a}{\sqrt5}\cdot\frac{2}{\sqrt5}=\frac{2a}{5}\).
因此 \(P\) 点到直线 \(L\) 的距离为 \(\frac{2a}{5}\).
\end{enumerate}
\solutionfigure{\bitmapfigure[max width=0.36\linewidth]{img/1959/national/q4_solution_fig1.png}}
\end{solution}

\begin{problem}
延长圆 \(O\) 的两弦 \(AB\)，\(CD\) 交于圆外一点 \(E\)，过 \(E\) 点作 \(DA\) 的平行线交 \(CB\) 的延长线于点 \(F\)，自 \(F\) 点作圆 \(O\) 的切线 \(FG\). 求证 \(FG=FE\).

\bitmapfigure[max width=0.42\linewidth,trim=0 0.45cm 0 0.15cm,clip]{img/1959/national/q5_fig1.png}
\end{problem}

\begin{solution}
由题意，\(A\)，\(B\)，\(E\) 共线，\(C\)，\(D\)，\(E\) 共线，且 \(C\)，\(B\)，\(F\) 共线. 因为 \(EF\parallel DA\)，所以 \(\angle FEB=\angle BAD\). 同弦 \(BD\) 所对的圆周角相等，故 \(\angle BAD=\angle BCD\). 又因为 \(CF\) 与 \(CB\) 在同一直线上、\(CE\) 与 \(CD\) 在同一直线上，所以 \(\angle BCD=\angle FCE\)，从而 \(\angle FEB=\angle FCE\). 同时，\(FB\) 与 \(FC\) 是同一条射线，所以 \(\angle EFB=\angle CFE\). 于是 \(\triangle EFB\sim\triangle CFE\). 对应边给出 \(\frac{FE}{FC}=\frac{FB}{FE}\)，从而 \(FE^2=FB\cdot FC\).

从点 \(F\) 引圆 \(O\) 的切线 \(FG\)，而直线 \(FBC\) 是圆 \(O\) 的割线. 根据切割线与切线定理，\(FG^2=FB\cdot FC\). 因而 \(FG^2=FE^2\). 因为 \(FG\)、\(FE\) 都是线段长度，均为正数，所以 \(FG=FE\).
\solutionfigure{\bitmapfigure[max width=0.42\linewidth,trim=0 0.45cm 0 0.15cm,clip]{img/1959/national/q5_solution_fig1.png}}
\end{solution}
